\section{問題定義}
\label{sec:problem}

\subsection{準備}

$\mathcal{I} = \{i_1, i_2, \ldots, i_m\}$ をアイテムの集合とする．
\emph{トランザクション} $T \subseteq \mathcal{I}$ はアイテムの部分集合である．
\emph{トランザクション系列} $\mathcal{D} = (T_1, T_2, \ldots, T_N)$ は，時間順に索引付けされた $N$ 個のトランザクションの順序付き集合である．
\emph{アイテムセット}（または\emph{パターン}）$P \subseteq \mathcal{I}$ はアイテムの空でない部分集合である．

\begin{definition}[ウィンドウサポート]
\label{def:windowed-support}
パターン $P$ とウィンドウサイズ $W \in \mathbb{N}^+$ に対して，位置 $t$ における\emph{ウィンドウサポート}は以下で定義される：
\begin{equation}
    \sigma_P(t) = \frac{|\{i \in [t, t+W) : P \subseteq T_i\}|}{W}, \quad t \in [1, N-W+1]
\end{equation}
得られる系列 $\mathbf{s}_P = (\sigma_P(1), \ldots, \sigma_P(N-W+1))$ を $P$ の\emph{サポート時系列}と呼ぶ．
\end{definition}

\subsection{因果帰属フレームワーク}

介入は既知の時刻 $t_0$ に発生するものとする．
目的は，\emph{処置パターン} $P^*$ のサポートに対するこの介入の因果効果を推定することである．

\begin{definition}[ドナープール]
\label{def:donor-pool}
処置パターン $P^*$ と候補パターン集合 $\mathcal{P}$ が与えられたとき，\emph{ドナープール}は以下で定義される：
\begin{equation}
    \mathcal{D}(P^*) = \{P \in \mathcal{P} : P \cap P^* = \emptyset\}
\end{equation}
ここで $P \cap P^*$ はアイテム集合の共通部分を表す．
ドナープール内のパターンは $P^*$ とアイテムを共有しない．
\end{definition}

アイテム非共有の根拠は，$P^*$ とアイテムを共有するパターンが同じ介入の直接的影響を受ける可能性があることにある（例：$i \in P^*$ に対するプロモーションは $i$ を含む全てのパターンに影響する）．

\begin{definition}[反事実サポート軌跡]
\label{def:counterfactual}
ドナープール $\mathcal{D}(P^*) = \{P_1, \ldots, P_J\}$ が与えられたとき，\emph{反事実サポート軌跡}は以下で定義される：
\begin{equation}
    \hat{\sigma}_{P^*}^{\mathrm{CF}}(t) = \sum_{j=1}^{J} w_j \cdot \sigma_{P_j}(t), \quad t \geq t_0
\end{equation}
ここで重み $\mathbf{w}^* = (w_1^*, \ldots, w_J^*)$ は介入前の予測誤差を最小化する：
\begin{equation}
    \mathbf{w}^* = \arg\min_{\mathbf{w} \in \Delta^J} \sum_{t=1}^{t_0-1} \left(\sigma_{P^*}(t) - \sum_{j=1}^{J} w_j \cdot \sigma_{P_j}(t)\right)^2
    \label{eq:weight-opt}
\end{equation}
ただし，単体制約 $\Delta^J = \{\mathbf{w} : w_j \geq 0, \sum_j w_j = 1\}$ に従う．
\end{definition}

\begin{definition}[サポートに対する因果効果]
\label{def:causal-effect}
時刻 $t \geq t_0$ における\emph{因果効果}は以下で定義される：
\begin{equation}
    \tau(t) = \sigma_{P^*}(t) - \hat{\sigma}_{P^*}^{\mathrm{CF}}(t)
\end{equation}
\emph{累積因果効果}は $\mathcal{T} = \sum_{t=t_0}^{T} \tau(t)$ である．
\end{definition}

\begin{definition}[アイテム非共有パターン制御]
\label{def:item-disjoint}
パターン $P_j$ が $P^*$ に対する妥当な\emph{アイテム非共有制御}であるとは，以下の条件を満たすときをいう：
\begin{enumerate}
    \item $P_j \cap P^* = \emptyset$（アイテム非共有），かつ
    \item $\sigma_{P_j}$ を $\sigma_{P^*}$ にフィッティングした際の介入前 RMSPE が閾値 $\delta$ 未満である．
\end{enumerate}
\end{definition}

\subsection{因果識別のための仮定}

本節では，因果効果の識別に必要な仮定を明示的に列挙し，パターンサポート文脈での妥当性を議論する．

\begin{definition}[SUTVA：安定単位処置値仮定]
\label{def:sutva}
各パターン $P_j \in \mathcal{D}(P^*)$ のサポート軌跡は，処置パターン $P^*$ に対する介入の有無に影響されない：
\begin{equation}
    \sigma_{P_j}^{\text{obs}}(t) = \sigma_{P_j}^{(0)}(t), \quad \forall j, \forall t
\end{equation}
\end{definition}

アイテム非共有基準（\Cref{def:item-disjoint}）は SUTVA の十分条件ではないが，その蓋然性を高める設計的保証である：
$P^*$ のアイテムを標的としたプロモーションは，$P^*$ とアイテムを共有しないパターンに直接的な影響を与えない．
ただし，間接的な波及効果（例：予算制約による代替購買）は SUTVA 違反となり得るため，\Cref{sec:experiments} で感度分析を行う．

\begin{definition}[線形因子モデル仮定]
\label{def:factor-model}
介入がない場合の全パターンのサポートが共通の潜在因子 $\boldsymbol{\mu}_t \in \mathbb{R}^r$ で駆動されるとする：
\begin{equation}
    \sigma_{P}^{(0)}(t) = \boldsymbol{\lambda}_P^\top \boldsymbol{\mu}_t + \epsilon_{P,t}, \quad \mathbb{E}[\epsilon_{P,t}] = 0
\end{equation}
ここで $\boldsymbol{\lambda}_P \in \mathbb{R}^r$ はパターン固有の因子負荷量，$r$ は因子数である．
\end{definition}

\begin{remark}
\Cref{def:factor-model} は平行トレンド仮定の一般化である．
DiD は $r=1, \boldsymbol{\mu}_t = t$ の特殊ケースに相当する．
SCM は，ドナープールの因子負荷量が処置パターンの因子負荷量を凸包内に含む場合に，反事実を無バイアスに復元できる~\cite{abadie2010synthetic}．
\end{remark}

\subsection{理論的保証}

\begin{theorem}[因子モデル下での一致性]
\label{thm:unbiased}
\Cref{def:factor-model} の線形因子モデルが成り立ち，SUTVA（\Cref{def:sutva}）が満たされ，かつ処置パターンの因子負荷量 $\boldsymbol{\lambda}_{P^*}$ がドナーの因子負荷量の凸包内にある（$\boldsymbol{\lambda}_{P^*} = \sum_j w_j^* \boldsymbol{\lambda}_{P_j}$，$\mathbf{w}^* \in \Delta^J$）とする．
このとき，介入前期間 $t_0$ が十分に長く，$\epsilon_{P,t}$ が弱従属で $\mathrm{Var}(\bar{\epsilon}) = O(1/t_0)$ を満たすならば，SCM 推定量は：
\begin{equation}
    |\hat{\tau}(t) - \tau^*(t)| = O_p\left(\frac{1}{\sqrt{t_0}}\right), \quad \forall t \geq t_0
\end{equation}
すなわち，介入前期間の長さに対して $\sqrt{t_0}$-一致性を持つ．
\end{theorem}

\begin{proof}
\Cref{eq:weight-opt} の最適化により $\hat{\mathbf{w}}$ が推定される．
因子モデルのもとで，介入前適合誤差は：
\begin{align}
    &\frac{1}{t_0}\sum_{t=1}^{t_0-1}\left(\sigma_{P^*}^{(0)}(t) - \sum_j \hat{w}_j \sigma_{P_j}(t)\right)^2 \\
    &= \frac{1}{t_0}\sum_{t=1}^{t_0-1}\left((\boldsymbol{\lambda}_{P^*} - \sum_j \hat{w}_j \boldsymbol{\lambda}_{P_j})^\top \boldsymbol{\mu}_t + \tilde{\epsilon}_t\right)^2
\end{align}
ここで $\tilde{\epsilon}_t = \epsilon_{P^*,t} - \sum_j \hat{w}_j \epsilon_{P_j,t}$ である．
凸包条件と介入前適合の最適性により $\hat{\mathbf{w}} \to \mathbf{w}^*$（$t_0 \to \infty$）が保証される．
すると介入後の予測誤差は $\hat{\sigma}^{\mathrm{CF}}(t) - \sigma_{P^*}^{(0)}(t) = (\hat{\mathbf{w}} - \mathbf{w}^*)^\top \boldsymbol{\lambda} \cdot \boldsymbol{\mu}_t + \tilde{\epsilon}_t = O_p(1/\sqrt{t_0})$ となる．
\end{proof}

\begin{remark}[有限標本バイアスと正則化]
\label{rem:regularization}
実践上，$t_0$ が小さい場合やドナー数 $J$ が大きい場合には SCM が過適合し，介入前適合が良好でも介入後バイアスが発生し得る．
この問題に対して，(1) $\ell_1$ ペナルティによるスパース重み推定，(2) リッジ正則化~\cite{amjad2018robust}，(3) Augmented SCM~\cite{benmichael2021augmented} による二重頑健推定が有効である．
本実装では，小規模ドナープール（$J \leq 20$）を前提とするため古典的 SCM を採用するが，大規模ドナープールへの拡張時には正則化が必須となる．
\end{remark}

\begin{theorem}[プラセボ検定の妥当性]
\label{thm:placebo}
帰無仮説 $H_0: \tau^*(t) = 0$（すべての $t \geq t_0$）のもとで，SUTVA が成立し，全 $J+1$ ユニット（処置 + ドナー）の介入後/介入前 RMSPE 比が交換可能であるならば，処置ユニットの比が最大となる確率は正確に $1/(J+1)$ である．
\end{theorem}

\begin{proof}
$H_0$ と SUTVA のもとでは，処置ユニットはドナーと統計的に区別できない．
交換可能性により，全 $(J+1)!$ 個の順位付けが等確率である．
特定のユニット（処置ユニットを含む）が最大の比を持つ確率は $1/(J+1)$ である．
\end{proof}

\begin{corollary}
$J$ 個のドナーによるプラセボ検定で達成可能な最小p値は $1/(J+1)$ である．
$p < 0.05$ を達成するには，少なくとも $J \geq 19$ 個のドナーが必要である．
\end{corollary}
